\documentclass[style=default]{genealogic}
\begin{document}

\tableofcontents

% --- RAMO MATERNO ---
\part{Ramo Materno: Famiglia Bianchi}
% Partiamo dalla bisnonna, quindi la prima generazione presente è la 1
\SetStartingGeneration{1}

\Generation % Visualizzerà: Generazione 1

\DeclarePerson{mario1880}[
  nome={Mario Rossi},
  padre={Gioacchino Rossi},
  nascita={1880},
  morte={1950}
]{
  Mario fu un pioniere della viticoltura nella valle. 
  La sua eredità agricola è documentata negli archivi comunali.
}

\Generation % Visualizzerà: Generazione 2
\DeclarePerson{giovanni1900}[
  nome={Giovanni Rossi},
  padre={\GetLinkName{mario1880}}, % Riferimento cliccabile
  nascita={1910}
]{
  Giovanni continuò il lavoro del padre, espandendo i possedimenti.
}

% --- RAMO PATERNO ---
\part{Ramo Paterno: Famiglia Esposito}
% Partiamo dal nonno, quindi la prima generazione presente è la 2
\SetStartingGeneration{2}

\Generation % Visualizzerà: Generazione 2 (perché Capitolo 1 + Offset 1)

\DeclarePerson{gennaro1885}[nome={Gennaro Esposito}]{Note}
\DeclarePerson{vincenzo1885}[nome={Vincenzo Esposito}]{Note}



\SetStartingGeneration{1}
\Generation % Deve stampare 1
\Generation % Deve stampare 2

\part{Altro Ramo}
\SetStartingGeneration{3}
\Generation % Deve stampare 3

%\Person{gennaro1885}{Gennaro Esposito}
%Il nonno paterno.

%\Generation % Visualizzerà: Generazione 3
%\Person{vincenzo1920}{Vincenzo Esposito}
%Tuo padre.

\end{document}